% =============================================================================
%  1-Introduccion.tex — Proemio, declaraciones y definiciones
% =============================================================================

Acuerdo interpersonal que celebran, por una parte, \textbf{\ParteA} (en adelante, \enquote{Rosario}) y, por la otra, \textbf{\ParteB} (en adelante, \enquote{Galileo}), a quienes en conjunto se les denominará \enquote{Las Partes}, al tenor de las siguientes declaraciones, definiciones y cláusulas.

\section{Declaraciones}

\begin{enumerate}[label=\Roman*., leftmargin=*, itemsep=0.4em]
  \item \textbf{Origen del acuerdo.} Rosario solicitó a Galileo que se le reconozca el derecho a:
  \begin{enumerate}[label=\alph*), itemsep=0.2em]
    \item Decir que no, expresando que eso es lo que necesita, sin estar obligada a dar explicaciones.
    \item Cerrar una discusión cuando ya no puede con ella, y detenerla en ese momento.
  \end{enumerate}

  \item \textbf{Aceptación condicionada.} Galileo acepta la solicitud, siempre y cuando se establezcan tres condiciones: (1) que quien plantee un tema exprese con claridad su opinión y la razón por la cual lo plantea; (2) que la pausa se solicite mediante frases clave claras y con reglas definidas; y (3) que existan plazos para retomar el tema que originó la pausa.

  \item \textbf{Bilateralidad.} Las Partes acuerdan que los derechos y obligaciones de este acuerdo aplican por igual a ambas: tanto de Rosario hacia Galileo como de Galileo hacia Rosario.

  \item \textbf{Motivo.} Las Partes reconocen que ciertos temas pueden afectarles sentimentalmente al grado de hacerles sentir vulnerados o invalidados, y que en esos momentos se necesita tiempo para reponerse antes de continuar.

  \item \textbf{Voluntad y alcance.} Las Partes celebran este acuerdo de manera libre, voluntaria y de buena fe, con el propósito de cuidar la relación y el bienestar emocional de ambas. Reconocen que no constituye un contrato con validez jurídica, sino un compromiso personal y moral entre ellas.
\end{enumerate}

\subsection*{Definiciones}

Para efectos de este acuerdo se entenderá por:

\begin{description}[leftmargin=0pt, itemsep=0.3em]
  \item[Parte solicitante:] la Parte que invoca la pausa activa.
  \item[Parte receptora:] la Parte a quien se le solicita la pausa activa.
  \item[Frase clave:] cualquiera de las frases establecidas en la Cláusula Cuarta, cuya sola mención detiene la conversación.
  \item[Pausa activa:] suspensión temporal de la conversación sobre un tema, sin que se suspenda la relación ni la interacción cotidiana entre Las Partes.
  \item[Estabilidad visible:] estado de regulación física, emocional o social de una Parte, perceptible por ambas Partes, que indica que los estímulos que la desbordaron han dejado de afectarla.
  \item[Suceso:] el momento en que se pronuncia la frase clave e inicia la pausa activa. Cada nueva pausa constituye un nuevo suceso.
  \item[Conversación de resolución:] conversación iniciada por cualquiera de Las Partes con el propósito de resolver, comprender o entender la causa que originó la pausa activa.
  \item[Conflicto solucionado:] estado que adquiere el tema que originó la pausa activa cuando no se retoma dentro del plazo establecido en la Cláusula Sexta.
\end{description}
